\chapter{Proofs for Huber's Further Robust Methods}
\label{annex:p15-huber-remaining-methods}

\begin{annexitem}{annex:p15.huber.mean.median}{thm:p15-huber-mean-median}{The first two source examples}
\begin{proof}
Writing \(\overline x=n^{-1}\sum x_i\), expansion gives
\(\sum(x_i-t)^2=\sum(x_i-\overline x)^2+n(t-\overline x)^2\).
For the absolute-value criterion, its left and right derivatives
are the numbers of observations strictly to the left minus those
on or to the right, and those on or to the left minus those strictly
to the right, respectively. A convex function is minimized exactly
when these derivatives bracket zero. This yields the stated
order-statistic interval, including ties.
The strong law gives consistency of the mean, and the i.i.d.\
central limit theorem gives its finite-variance conclusion.

For the median, the local CDF expansion gives
\(F(-\delta)<1/2<F(\delta)\) for all sufficiently small positive
\(\delta\). The strong law for the two indicator counts puts
both central order statistics between \(-\delta,\delta\) eventually.
For a fixed \(u\), put \(t_n=u/\sqrt n\).
The count \(N_n=\sum_i{\bf1}_{\{X_i\leq t_n\}}\) is binomial with
\[
 p_n=F(t_n)=1/2+f_0u/\sqrt n+o(n^{-1/2}).
\]
Its centered summands are bounded and its variance is asymptotic
to \(n/4\), so the Lindeberg theorem applies. The event
\(X_{(j)}\leq t_n\) is \(N_n\geq j\). For either central index
\(j=n/2+O(1)\), the normalized threshold tends to \(-2f_0u\).
Thus each central order statistic has limiting CDF \(\Phi(2f_0u)\).
The midpoint lies between them; their two probability bounds give
the same limit. Substituting \(\phi(0)=1/\sqrt{2\pi}\) gives
variance \(\pi/2\). Neither the indicator law nor this argument
requires a second moment of \(X\).
\end{proof}
\end{annexitem}

\begin{annexitem}{annex:p15.huber.scale.model}{thm:p15-huber-scale-model}{The source's one-sided scale branch and its boundary qualification}
\begin{proof}
For \(X\sim N(0,1)\),
\(G_Y(u)=2\Phi(e^{u/2})-1\). Differentiation gives
\(g_Y(u)=e^{u/2}\phi(e^{u/2})\) and the displayed logarithmic
derivative. Its derivative is \(e^u/2>0\), proving log-concavity.
The equation for either clipping tangent is
\((e^u-1)/2=\pm c\), giving the two stated regimes.
The general positive-contamination construction, including its
normalization and ordered minimax proof, therefore applies directly.

For \(c\geq1/2\) the lower tangent is infinite, and
\(G_Y(\log q^2)=A(q)\), \(g_Y(\log q^2)=q\phi(q)\).
Substituting \(c=(q^2-1)/2\) proves the source formula.
The least-favorable excess law is entirely on the upper clipped
tail. Since its score mean is zero, its mixture representation
gives
\((1-\varepsilon)E_{G_Y}\chi_c=-\varepsilon c\).
Thus every centered competitor must have \(E_H\chi_c=c\).
The score never exceeds \(c\), attaining it exactly when
\(|X|\geq q\); this proves the support condition.
Equality at the boundary cannot be discarded when stating what
centering itself implies.

If there is no mass at \(|X|=q\), bounded difference quotients
for the piecewise smooth Lipschitz transformed score permit
differentiation of its population mean by dominated convergence.
Its derivative with respect to its argument is
\(e^u/2\) below the upper join and zero above it.
Every strict-tail contaminant therefore contributes zero to
that derivative. The remaining slope is
\[
 \frac{1-\varepsilon}{2}E_\Phi[X^2{\bf1}_{\{|X|<q\}}]
 =\frac{1-\varepsilon}{2}\{A(q)-2q\phi(q)\}=I_Y.
\]
The efficient-score information identity from the prototype theorem
shows that the centered score second moment is also \(I_Y\).
The convex normality theorem gives variance \(1/I_Y\).
Scale transformation adds \(\log\sigma^2\) to every transformed
observation, proving the asserted equivariance and normalization.
The general prototype minimax risk proof applies to this same
centered class and parameter \(\log\sigma^2\).

At a boundary atom the transformed score's left derivative is
\(q^2/2\), and its right derivative is zero. There is no symmetry
of the transformed law requiring these slopes to average as in
the location problem. The two population slopes can differ.
The active one-sided convex-score limit therefore applies when
its premises hold; a single normal limit has not followed from
centering alone.
\end{proof}
\end{annexitem}

\begin{annexitem}{annex:p15.huber.scale.computation}{thm:p15-huber-scale-computation}{Scale equations, zero observations and a worked sample}
\begin{proof}
For a nonzero observation,
\[
 \chi_c(\log x^2-\log S^2)
 =\tfrac12\bigl(\min\{u,\max(\ell,x^2/S^2)\}-1\bigr).
\]
When \(c\geq1/2\), the lower nominal score limit is \(-1/2\),
which explains \(\ell=0\). When \(c<1/2\), it is clipped at
\(-c\), giving \(\ell=1-2c\).
The right side has a well-defined value also for \(x=0,S>0\).
This algebraic extension does not claim that \(\log0\) is a
real-valued observation in the transformed model.

Each summand is continuous and nonincreasing in positive \(S\).
Its two limits give those in the statement. If the upper limit
equals \(n\), all nonzero observations are saturated at \(u\)
for every sufficiently small positive \(S\), while all zero
observations equal \(\ell\); hence an interval of exact roots
abuts zero. If the upper limit is strictly larger than \(n\),
the intermediate value theorem and the strict limiting inequalities
put the roots in a compact positive interval. Monotonicity makes
the zero set an interval. For \(\ell=0\), every nonzero observation
that is not upper-saturated contributes a strictly decreasing
positive term. At a root below the upper limit there must be
one such term, proving uniqueness.
When \(\ell>0\), every summand can simultaneously be at one
of its two caps, producing a plateau at the target; half of the
observations at each cap is one way this occurs.
Midpoint selection for \(\log S^2\) exponentiates to the geometric
mean of the two positive scale endpoints.

For the five-point sample at \(q=2,S=2\), four terms are \(1/4\)
and the outlier term is four, summing to five. The strict upper
limit ensures uniqueness. For four observations with only one
nonzero value one, the upper limit is exactly four and all
\(0<S\leq1/2\) keep that observation saturated.
For five observations with only that nonzero value, the upper
limit four is below the target five. These give the claimed
nonuniqueness and nonexistence examples.
\end{proof}
\end{annexitem}

\begin{annexitem}{annex:p15.huber.scale.gauging}{thm:p15-huber-scale-gauging}{Normal calibration and historical Table III}
\begin{proof}
Under the normal law, the population equation is
\[
 E_\Phi\min(X^2/S^2,q^2)=1.
\]
Its left side is continuous, strictly decreasing in \(S>0\),
with limits \(q^2>1\) and zero. Strictness follows from the
positive normal mass of unsaturated nonzero observations.
At \(S=1\) it is \(B(q)<EX^2=1\), so the unique root is below one.
Multiplying the equation by \(S^2\) gives \(B(qS)=S^2\).
The strong law at any fixed scales on opposite sides of this
root, together with monotonicity, proves convergence of the
sample root to \(S_\infty\).

Write \(S_n=S_\infty S_n'\). The clipping identity
\[
 \min\{X^2/(S_\infty^2S_n'^2),q^2\}
 =S_\infty^{-2}\min\{X^2/S_n'^2,(qS_\infty)^2\}
\]
converts the sample equation to the stated normal-calibrated one.
The calibration relations follow. Equivalently
\(B(q')/q'^2=1/q^2\); its left side decreases strictly from
one to zero, by the same positive-mass argument, establishing
uniqueness of \(q'\).
In the earlier calibrated contamination class \(S_n\to1\),
so dividing by \(S_\infty\) gives limit \(1/S_\infty\), not one.
The numerical example evaluates these exact equations at \(q=2\).
The root-mean-square comparison follows by summing the five
squared observations.
\end{proof}
\end{annexitem}

\begin{annexitem}{annex:p15.huber.trimming.calculation}{thm:p15-huber-trimming-calculation}{Global computation, deletion and a failed consistency inference}
\begin{proof}
The criterion is continuous. Outside
\([\min_i x_i-k,\max_i x_i+k]\) it equals \(nk^2/2\), whereas at
any observation its value is at most \((n-1)k^2/2\).
It therefore attains a global minimum on that compact interval.
Between consecutive breakpoints \(x_i\pm k\), the active set is fixed
and the criterion is a quadratic with derivative
\(\sum_{i\in A}(t-x_i)\). If \(A\ne\varnothing\), minimize this
quadratic on the closed interval by its mean, clipped to the
interval endpoints. If \(A=\varnothing\), the criterion is constant.
Comparing the finitely many candidate values and all breakpoints
proves the algorithm. At a minimum without a boundary residual
the active set cannot be empty, since the value would be \(nk^2/2\);
the derivative condition gives its arithmetic mean.

Suppose \(r\) removed observations all have residual magnitude
greater than \(k\) at a global minimizer \(t_0\). For every \(t\),
\[
 Q_{\rm retained}(t)
 =Q_n^{\rm tr}(t)-Q_{\rm removed}(t)
 \geq Q_n^{\rm tr}(t_0)-rk^2/2
 =Q_{\rm retained}(t_0).
\]
This is a global argument; mere stationarity would not suffice.
For \((0,1,2,100)\), on \(0\leq t\leq1\) the criterion is
\(t^2-t+3/2\); on \(1\leq t\leq2\) it is
\((t-3/2)^2+5/4\). All remaining breakpoint intervals have
minimum at least \(3/2\): outside the three-point cluster only
one of its observations can contribute below \(1/2\), or the
remote observation alone can do so. The two stated minimizers
and the value at their midpoint follow.
At \(t=3/2\), the clipped-linear score contributions are
\(-1,-1/2,1/2,1\), summing to zero; the two middle contributions
make that monotone score root unique. The mean and median are
direct arithmetic.

For the population derivative write
\[
 \lambda(t)=\int_{t-k}^{t+k}(x-t)f(x)\,dx.
\]
The change in the interior factor contributes
\(-\int_{-k}^kf\). Continuity of \(f\) at the two moving boundaries
makes their contributions \(k f(k)+k f(-k)\).
This proves the displayed derivative, and identifies the formal
variance denominator used in the source. It is not a proof of
the asserted normal limit for this discontinuous score.

For the stipulated mixture, the two remote intervals are outside
\([-1,1]\), the central interval has mass \(1/10\), and the density
is zero near \(\pm1\). Thus \(D=1/10\). The centered uniforms
have variance \(1/300\). The expected criterion at zero is
\[
 (1/10)(1/600)+(9/10)(1/2)=2701/6000,
\]
while its value at \(10\) is
\[
 (9/20)(1/600)+(11/20)(1/2)=1103/4000,
\]
which is strictly smaller.
If a global empirical minimizer converged in probability to zero,
the deterministic Lipschitz bound
\[
 |n^{-1}Q_n^{\rm tr}(T_n)-n^{-1}Q_n^{\rm tr}(0)|\leq |T_n|
\]
and the strong law at the fixed points \(0,10\) would contradict
its minimizing property. Hence positivity of the local derivative
quantity has not supplied global consistency.
\end{proof}
\end{annexitem}

\begin{annexitem}{annex:p15.huber.smooth.stationary.limit}{thm:p15-huber-smooth-stationary-limit}{The separate smooth nonconvex score argument}
\begin{proof}
Uniform continuity of \(\psi'\) gives a modulus
\(\omega(s)=\sup_{|x-y|\leq s}|\psi'(x)-\psi'(y)|\) finite for
finite \(s\) and tending to zero at zero. Finiteness at larger
\(s\) follows by dividing an interval into finitely many smaller ones.
The fundamental theorem of calculus gives the exact identity
\[
 0=\frac1n\sum_i\psi(X_i)-T_n D_n,\qquad
 D_n=\frac1n\sum_i\int_0^1\psi'(X_i-vT_n)\,dv .
\]
The difference between \(D_n\) and \(n^{-1}\sum_i\psi'(X_i)\)
is at most \(\omega(|T_n|)\), which tends to zero in probability.
The strong law and the absolute-integrability premise give
\(D_n\to d>0\) in probability.
Since \(T_n\to0\), the identity forces \(n^{-1}\sum_i\psi(X_i)\to0\)
in probability. Its strong-law limit is \(E\psi(X)\), so that
expectation is zero. With probability tending to one the denominator
is nonzero, and
\[
 \sqrt n T_n=
       \frac{n^{-1/2}\sum_i\psi(X_i)}{D_n}.
\]
The i.i.d.\ central limit theorem and the continuous perturbation
rule prove the claim, including a zero numerator variance.
The proof assumes a consistent root; it does not infer one from
the local equation.
\end{proof}
\end{annexitem}

\begin{annexitem}{annex:p15.huber.exact.median.bias}{thm:p15-huber-exact-median-bias}{Contamination bias and indistinguishable locations}
\begin{proof}
Let \(a=1-\varepsilon>1/2\). For every contaminant,
\(a\Phi(x)\leq F(x)\leq a\Phi(x)+\varepsilon\).
If \(x>c_\varepsilon\), \(F(x-)\geq a\Phi(x)>1/2\);
if \(x<-c_\varepsilon\), \(F(x)\leq a\Phi(x)+\varepsilon<1/2\).
These exclude every median outside the asserted interval.
Putting all contaminant mass at \(M>c_\varepsilon\) makes
\(c_\varepsilon\) the unique median; reflection attains the other end.

For the sharp lower bound form the even density
\[
 p(x)=a\max\{\phi(x-c_\varepsilon),\phi(x+c_\varepsilon)\}
     =a\phi(|x|-c_\varepsilon).
\]
Its integral is \(2a\Phi(c_\varepsilon)=1\).
It dominates the normal component with center \(c_\varepsilon\)
and also the one with center \(-c_\varepsilon\).
Subtracting either component leaves a nonnegative density of
mass \(\varepsilon\), so the same observed law is in both
location-contamination models. Any reported real location \(t\)
has \(\max(|t-c_\varepsilon|,|t+c_\varepsilon|)\geq c_\varepsilon\).
For a translation-equivariant functional these two relative errors
translate into errors on the two corresponding zero-centered
contamination laws. Thus no such functional improves the bound,
while the median attains it.
When \(a\leq1/2\), the integral of
\(a\max\{\phi(x-c),\phi(x+c)\}\) is \(2a\Phi(c)\leq1\)
for every finite \(c\). Add any nonnegative residual density of
the missing mass. The same indistinguishability argument for
arbitrarily large \(c\) proves the infinite lower bound.
The median's error is also unbounded, by contaminants at an
arbitrarily remote point.

For the source's flat construction let
\[
 p_{\rm flat}(x)=
 \begin{cases}
 a\phi(0),&|x|<c_{\rm flat},\\
 a\phi(|x|-c_{\rm flat}),&|x|\geq c_{\rm flat}.
 \end{cases}
\]
Its integral is \(a(1+2c_{\rm flat}\phi(0))=1\).
It dominates both shifted normal components, so it does give
a valid indistinguishability lower bound.
But
\[
 \frac{\varepsilon}{2a}
 =\Phi(c_\varepsilon)-\tfrac12
 =\int_0^{c_\varepsilon}\phi(t)\,dt
 <c_\varepsilon\phi(0)
\]
for \(\varepsilon>0\). Hence its bound is strictly smaller than
the exact median bound. Differentiability of \(\Phi\) at zero gives
the first-order equality as \(\varepsilon\downarrow0\).
The stated decimal comparisons result by evaluating the exact
inverse-CDF and flat-bound formulas, not by treating the approximation
as a finite-\(\varepsilon\) inequality in the other direction.
\end{proof}
\end{annexitem}

\begin{annexitem}{annex:p15.huber.cdf.domain}{thm:p15-huber-cdf-domain}{Compactness and the dense finite-information class}
\begin{proof}
For symmetric \(F\), and \(t>0\), put \(R_F(t)=F([-t,t])\).
Direct splitting at zero gives
\[
 B_F(t)=\tfrac12+\tfrac12R_F(t),\qquad
 B_F(t-)=\tfrac12+\tfrac12F((-t,t)).
\]
These identities include an atom at zero. The analogous negative
identities follow by symmetry. Since \(\Phi\) is continuous, a
CDF bound holds also for the left limits.
If symmetric subprobabilities converge vaguely, their masses
on \([-t,t]\) converge at every \(t\) for which the limit gives
no mass to \(\{-t,t\}\). To verify this interval assertion,
sandwich its indicator by compact continuous ramps with supports
slightly larger and smaller; their errors in the limit tend to
zero by the absence of boundary atoms.
At all such \(t\) the required inequalities pass to the limit.
Right and left limits then give them at every remaining point,
including zero. Symmetry itself is closed by its \(C_c\) integral
identities. Thus \(\mathcal K_\varepsilon\) is vaguely closed in
the compact subprobability domain of
\cref{thm:p15-vague-subprobability-compactness}.
It is convex because \(B_F\) depends affinely on \(F\).
Letting \(x\to\infty\) in the bound gives
\(|(1+m)/2-1|\leq\varepsilon\), or \(m\geq1-2\varepsilon\).

Here is an approximation that stays inside the neighborhood.
For \(0<\delta<1\), put \(F_\delta=(1-\delta)F+\delta\Phi\);
its completed CDF differs from \(\Phi\) by at most
\((1-\delta)\varepsilon\).
Convolve its finite real-line measure with \(N(0,s^2)\).
The completed CDF of the resulting measure is
\(E B_{F_\delta}(x-sZ)\), \(Z\sim N(0,1)\); the equal missing
masses remain at the two infinities. Since the normal CDF is
Lipschitz with constant \(\phi(0)\),
\[
 \sup_x|E\Phi(x-sZ)-\Phi(x)|
       \leq s\phi(0)E|Z|=s/\pi .
\]
Choosing \(s\leq\pi\delta\varepsilon\) keeps the convolved law
inside \(\mathcal K_\varepsilon\).

It has a smooth positive density \(f_s(x)=\int\phi_s(x-y)\,F_\delta(dy)\).
Differentiation is justified by bounded Gaussian derivatives.
Pointwise Cauchy--Schwarz gives
\[
 \frac{f_s'(x)^2}{f_s(x)}
 \leq\int\frac{(x-y)^2}{s^4}\phi_s(x-y)\,F_\delta(dy).
\]
Tonelli and the Gaussian second moment bound its information by
\(F_\delta(\mathbb R)/s^2<\infty\).
The finite-information characterization therefore applies.
As \(\delta\downarrow0\) and \(s\downarrow0\), these measures
converge vaguely to \(F\): a compactly supported continuous test
is bounded and uniformly continuous, so Gaussian smoothing
converges uniformly on its translates, by splitting small and
large Gaussian displacements. The additional mixture tends to
the original measure as well. This proves the required density.
\end{proof}
\end{annexitem}

\begin{annexitem}{annex:p15.huber.cdf.construction}{thm:p15-huber-cdf-construction}{The cosine core, normal middle and exponential tail}
\begin{proof}
The map \(z\mapsto2z\tan z\) increases continuously from zero
to infinity on \((0,\pi/2)\), proving existence and uniqueness of
\(z\). It gives \(c\tan(ca/2)=a\).
Moreover \(c^2=2z/\tan z\) decreases from two to zero as \(a\)
increases from zero to infinity. To verify the monotonicity,
\(\tan z/z\) has derivative with numerator
\(z\sec^2z-\tan z\); this numerator starts at zero and has
positive derivative \(2z\sec^2z\tan z\).

Let \(f_a(x)=\phi(a)\cos^2(cx/2)/\cos^2(ca/2)\), \(0\leq x\leq a\).
For \(0<x<a\), the same monotonicity of \(\tan u/u\) gives
\(c\tan(cx/2)<x\). Hence
\[
 \frac{d}{dx}\log(f_a(x)/\phi(x))
       =x-c\tan(cx/2)>0.
\]
The ratio equals one at \(a\); thus \(f_a<\phi\) inside, proving
\(E(a)>0\). The integral of \(\cos^2(cx/2)\) gives exactly the
formula defining \(E(a)\).
Implicit differentiation gives \(c'(a)<0\), and differentiation
of the logarithm of \(f_a\), using \(c\tan(ca/2)=a\), yields
\[
 \partial_a\log f_a(x)
   =c'(a)\{a\tan(ca/2)-x\tan(cx/2)\}<0\quad(0\leq x<a).
\]
The moving-endpoint contribution to
\(\int_0^a(\phi-f_a)\) is zero. Differentiation under the integral
on any compact positive \(a\)-interval is justified by continuity
of these functions on the corresponding compact triangle.
It follows that \(E'(a)>0\).
As \(a\downarrow0\), \(0<E(a)<\Phi(a)-1/2\to0\).
For large \(a\), \(z\geq\pi/4\) and
\[
 \cos^{-2}z=1+\frac{a^4}{4z^2}
       \leq1+4a^4/\pi^2.
\]
Therefore \(\int_0^af_a\leq a\phi(a)(1+4a^4/\pi^2)\to0\),
so \(E(a)\to1/2\).

The tail equation uses
\[
 L(b)=\phi(b)/b-(1-\Phi(b))
     =\int_b^\infty(t/b-1)\phi(t)\,dt>0 .
\]
It has derivative \(-\phi(b)/b^2<0\), limit infinity at zero,
and limit zero at infinity. Thus its root for every positive
\(\varepsilon\) is unique. As \(\varepsilon\) increases in
\((0,1/2)\), \(a(\varepsilon)\) strictly increases from zero
to infinity, while \(b(\varepsilon)\) strictly decreases from
infinity to a finite positive limit. Their unique crossing
defines \(\varepsilon_*\). This is the exact range condition;
the source's approximate decimal is not used as a proof of it.

On the permitted range \(a\leq b\), the central mass on the
positive half-line is \(\Phi(a)-1/2-\varepsilon\), the middle
mass is \(\Phi(b)-\Phi(a)\), and the tail mass is
\(\phi(b)/b=1-\Phi(b)+\varepsilon\).
Their sum is \(1/2\), proving normalization.
Values match at both joins; their log derivatives match because
the core score at \(a\) is \(a\), and the middle score at \(b\)
is \(b\). The even extension is \(C^1\) at zero as well.
The candidate is strictly positive.

On \(0<x<a\), \(f_a<\phi\), so the candidate CDF difference
decreases from zero to \(-\varepsilon\). It is constant through
the normal middle. On \(x>b\), its tail-to-normal density ratio
is \(\exp((x-b)^2/2)\geq1\), so the CDF difference increases back
to zero. Symmetry gives the negative half-line bound. This proves
membership in the neighborhood, including the crossing \(a=b\).

Its score is the displayed bounded continuous piecewise smooth
function. Integrating \(\psi_*^2 f_*\) over the core uses
\(\tan^2(cx/2)\cos^2(cx/2)=\sin^2(cx/2)\), whose integral is
\(a/2-\sin(ca)/(2c)\). Over the middle,
\(\int_a^b x^2\phi(x)\,dx=\Phi(b)-\Phi(a)-b\phi(b)+a\phi(a)\).
Over the two tails the information is \(2b\phi(b)\).
Adding these terms proves the exact information formula.
As \(\varepsilon\downarrow0\), the core width tends to zero
and the tail cutoff tends to infinity; the density and score
reduce to the standard normal construction.
\end{proof}
\end{annexitem}

\begin{annexitem}{annex:p15.huber.cdf.saddlepoint}{thm:p15-huber-cdf-saddlepoint}{The source variational proof and the atomic boundary cases}
\begin{proof}
Put \(w(x)=2\chi_*(x)-\psi_*(x)^2\), an even bounded function.
For \(r=|x|\) away from the joins it is
\[
 w(r)=c^2\ (r<a),\qquad
       2-r^2\ (a<r<b),\qquad -b^2\ (r>b).
\]
The chosen derivative representative makes its value at each join
the average of its actual one-sided limits.
The number \(d_a=c^2+a^2-2\) is nonnegative.
Indeed, on \((0,a)\) the derivative of the core score is increasing,
and its integral is \(\psi_*(a)-\psi_*(0)=a\).
Its endpoint value \((c^2+a^2)/2\) is therefore at least its
average one.

For \(t>0\) define
\[
 \overline R_F(t)=F(\{|x|<t\})+\tfrac12F(\{|x|=t\}),\qquad
 D(t)=F([-t,t])-F_*([-t,t]),
\]
and let \(\overline D(t)=\overline R_F(t)-F_*([-t,t])\).
The candidate CDF is \(\Phi(t)-\varepsilon\) on \([a,b]\).
The completed CDF constraints and their left limits, together
with the radial identities proved above, imply that \(D(t)\)
and \(\overline D(t)\) are nonnegative there.
Express \(w\) as the constant \(-b^2\), plus a downward-step
contribution of size \(d_a\) at \(a\), one of size two at \(b\),
and a continuous decrease with derivative \(-2t\) between them.
More explicitly use indicators \({\bf1}_{\{r<t\}}\), with half
weight at equality for the two discrete steps, and
\(\int_a^b2t{\bf1}_{\{r<t\}}\,dt\).
Integrating this identity against \(F-F_*\) gives
\[
 \int w\,d(F-F_*)=
 b^2(1-m)+d_a\overline D(a)+2\overline D(b)
                  +\int_a^b2tD(t)\,dt\ \geq0.
\]
The value of the integral at countably many radial atoms does not
affect its Lebesgue integral. If \(a=b\), the two step sizes add
to \(c^2+b^2\) at the same point, precisely the actual single jump.
This establishes the source's cumulative-function integration
argument with its mass term and all boundary-atom weights explicit.

Since \(f_*\) is a positive \(C^1\) density with bounded score,
integration by parts gives
\(\int\chi_* f_*=\int\psi_*^2f_*=I_*\); thus
\(\int w\,dF_*=I_*\).
The preceding inequality proves \(2A_F-B_F\geq I_*\) for
every member, without assuming it has a density.
For finite-information members, the weak-derivative integration
identity in \cref{thm:p15-information-saddlepoint-criterion}
and Cauchy--Schwarz give
\[
 I(F)\geq A_F^2/B_F\geq2A_F-B_F\geq I_*.
\]
Zero denominators are interpreted as in the stated score domain;
the positive last bound forces \(A_F>0\) if \(B_F=0\).
Infinite-information members cannot improve the minimum.
The equality-components theorem gives uniqueness because \(f_*\)
is positive everywhere. At that density every admissible opposing
score has payoff at most \(I_*\); hence this is a saddlepoint.

For a proper symmetric law, the odd continuous bounded score has
mean zero. It is globally Lipschitz. Its difference quotients are
therefore uniformly bounded, and their one-sided pointwise limits
are its one-sided derivatives. At a positive join and its negative
reflection, the left and right derivatives are interchanged.
Symmetry gives equal masses to those points, so their integrated
left and right derivatives agree and equal the average convention.
Dominated convergence consequently gives the actual derivative
\[
 \lambda_F'(0)=-\int\chi_*\,dF=-A_F<0.
\]
This also proves the statement at the merged join \(a=b\).
The second score moments are continuous under small translations
by bounded convergence. The score's monotonicity and negative
population derivative give the sign and consistency premises;
the convex primitive is coercive because its tail slopes are
\(-b,b\). The active convex normality theorem yields variance
\(B_F/A_F^2\leq1/I_*\).
Here \(B_F>0\): otherwise the proper law is concentrated at the
only zero of the strictly increasing central score, namely zero,
whose CDF distance from the standard normal is \(1/2\).
That law is outside this neighborhood.

At \(F_*\) equality holds. Bounded continuous truncated squared
losses converge under the normal limit; sending the truncation
level to infinity gives the asserted pointwise limiting risk
and then its supremum. For any translation-equivariant estimator,
\cref{thm:p15-location-integrated-score-bound} at \(f_*\) gives
the finite-\(C\) lower bound
\((I_*+4\pi^2/C^2)^{-1}\). Its ordered limits prove the
all-equivariant lower bound without interchanging a supremum
over laws and a sample-size limit.
\end{proof}
\end{annexitem}

\begin{annexitem}{annex:p15.huber.joint.finite}{thm:p15-huber-joint-finite}{The exact existence condition, finite computation and iteration}
\begin{proof}
Fix a tentative location \(t\). The function
\[
 L_t(s)=\sum_i\min((x_i-t)^2/s^2,k^2),\qquad s>0,
\]
is continuous and nonincreasing, with limits
\(k^2(n-\#\{i:x_i=t\})>n\beta\) and zero.
At level \(n\beta\), at least one nonzero residual is unsaturated:
otherwise the value would still be its upper limiting value.
Such a summand strictly decreases. Hence the equation
\(L_t(s)=n\beta\) has exactly one solution \(s_t>0\).
Continuity of \(s_t\) follows by choosing scales just below and
above its root, preserving their strict inequalities for all
locations close to \(t\), and shrinking that bracket.
If \(M=\max_i|x_i|\) and \(|t|>M\), comparison of the residuals
with \(|t|-M\) and \(|t|+M\) gives
\[
 (|t|-M)/\sqrt\beta\leq s_t\leq(|t|+M)/\sqrt\beta .
\]
Thus \(s_t/|t|\to1/\sqrt\beta\) and the average first score tends
to \(-\sqrt\beta\) as \(t\to\infty\), and to \(+\sqrt\beta\)
as \(t\to-\infty\). Its continuity and the intermediate value
theorem supply a joint solution.

At any solution, if \(m_2\) observations are strictly interior,
\((n-m_2)k^2\leq n\beta\), so
\(m_2\geq n(1-\beta/k^2)\). The multiplicity hypothesis ensures
that these interior observations are not all equal.
For uniqueness suppose \((T_1,S_1)\) and \((T_2,S_2)\) solve
the equations. If the scales agree, the first score is strictly
decreasing near a root with an interior observation and globally
nonincreasing, so the locations agree.
Otherwise take \(S_1<S_2\), put
\(d=(T_2-T_1)/(S_2-S_1)\), \(T(s)=T_1+d(s-S_1)\), and
\[
 R(s)=\tfrac12\sum_i\{\psi(k,(x_i-T(s))/s)^2-\beta\}
               +d\sum_i\psi(k,(x_i-T(s))/s).
\]
This function is absolutely continuous on the compact positive
scale interval. Away from its finitely many joining points,
direct differentiation gives
\[
 R'(s)=-s^{-1}\sum_{|(x_i-T(s))/s|<k}
                       (d+(x_i-T(s))/s)^2\leq0.
\]
Both endpoint values are zero. Therefore the derivative vanishes
almost everywhere. Interior observations at \(S_1\) remain interior
on a short interval to its right; each of their squared terms
must vanish there. This forces all of them to equal
\(T_1-dS_1\), contradicting their distinctness.
This is the source's monotonicity proof with the piecewise derivative
and the missing positive-\(\beta\) premise made explicit.

In a fixed active partition, multiply the two score equations by
\(S\) and \(S^2\). They become
\[
 \sum_{\rm interior}x_i-m_2T+(m_3-m_1)kS=0,\qquad
 \sum_{\rm interior}(x_i-T)^2+(m_1+m_3)k^2S^2=n\beta S^2.
\]
Eliminating \(T\) gives the displayed formulas.
The interior sum of squares about its mean is positive because
the interior values are not all equal; the denominator is therefore
positive as well. Every true solution appears in the finite
enumeration of ordered active blocks, and a candidate passing
its interval conditions solves the original equations. Uniqueness
proves the algorithm.
The source's iteration uses the same formulas, but the next
candidate depends only on the current finite partition. If every
update is defined, the sequence after its first step has only
finitely many possible states. Repetition gives a period; a fixed
point is a solution. Neither the finite state count nor uniqueness
excludes a period greater than one.

For the worked example the three interior observations have mean
zero and squared-deviation sum two. The denominator is
\(5/2-2=1/2\), giving \(S=2,T=0\), and \(M>2\) verifies the
partition. If all observations equal \(x\), the first equation
forces \(T=x\); the second then reads \(0=n\beta\).
It fails for positive \(\beta\) and leaves every \(S>0\) possible
when \(\beta=0\).
\end{proof}
\end{annexitem}

\begin{annexitem}{annex:p15.huber.joint.regular.limit}{thm:p15-huber-joint-regular-limit}{A full regular-domain justification of the location variance}
\begin{proof}
Write \(r=\beta/k^2\). The normal component makes the mass of
every atom of \(F\) at most \(\varepsilon<1-r\).
We first justify the eventual sample condition. On an event of
probability one, the empirical masses of intervals with endpoints
in a fixed dense set of non-atoms converge to their true masses,
by the strong law. The empirical laws are tight there, by choosing
large such intervals. If a subsequence of empirical atoms had
mass greater than \(\varepsilon+\delta\), tightness would keep
their locations bounded; a further subsequence converges to \(t\).
Intervals shrinking to \(t\), with non-atomic endpoints, would then
force \(F\{t\}\geq\varepsilon+\delta\), a contradiction.
Thus the maximal sample multiplicity divided by \(n\) has limsup
at most \(\varepsilon\), proving eventual existence and uniqueness
by the finite theorem.

For the population equations, at every fixed \(t\) the second
score moment as a function of \(s>0\) decreases continuously from
\(k^2(1-F\{t\})>\beta\) to zero. The positive normal density
makes its decrease strict at a root. It has a unique continuous
root \(s_t\). Tightness of \(F\), or bounded convergence along
subsequences of \(s_t/|t|\), gives \(s_t/|t|\to1/\sqrt\beta\)
as \(|t|\to\infty\): ratios tending to zero would give moment
\(k^2\), ratios tending to infinity would give zero, and a finite
positive limit \(v\) must satisfy \(\min(v^{-2},k^2)=\beta\).
Consequently the first population score crosses zero.
The uniqueness proof with the function \(R(s)\) above applies
after replacing sums by expectations. Its differentiation is
justified by bounded derivatives on a compact positive scale
interval. Vanishing of its derivative would require the normal
component within a nonempty interior interval to concentrate at
one point, impossible. At equal scales the first score is strictly
decreasing because of that component. Hence the population pair
is unique. Symmetry makes its location zero, giving the stated \(S_0\).

We supply the compactness step needed for consistency.
On the almost-sure interval-convergence event above, consider
any subsequence of sample roots. If \(|T_n|\) tends to infinity
while \(S_n\) stays bounded, empirical tightness makes the first
score average tend to \(k\) or \(-k\), a contradiction.
If \(S_n\to\infty\), take a subsequential limit of \(T_n/S_n\)
in the extended line. Empirical tightness makes the two moment
averages tend to \(\psi(k,-v)\) and its square for a finite
limit \(v\), or to \(\pm k,k^2\) for an infinite one.
The equations require zero and \(\beta>0\), again impossible.
This bounds locations and scales above.
If \(S_n\downarrow0\), a bounded location subsequence tends to \(t\).
Every observation outside a fixed small neighborhood of \(t\)
is then saturated. Interval convergence followed by shrinking the
neighborhood forces
\(\beta\geq k^2(1-F\{t\})\geq k^2(1-\varepsilon)>\beta\),
a contradiction. The roots therefore remain in a compact
location-scale set with positive lower scale, pathwise.

On such a compact parameter set both score functions
\(\psi(k,(x-t)/s)\) and its square are bounded and uniformly
Lipschitz in \((t,s)\), uniformly in \(x\): derivatives vanish
outside the clipped region, and within it the standardized
residual is bounded by \(k\). A finite parameter grid and the
strong law prove uniform convergence of both empirical moment
functions. Every accumulation point of their roots solves the
population equations; uniqueness proves almost-sure consistency.

For the limit distribution let
\[
 m(x;t,s)=\bigl(\psi(k,(x-t)/s),\,
                      \psi(k,(x-t)/s)^2-\beta\bigr)^{\mathsf T}.
\]
Away from boundary residuals its Jacobian is
\[
 -s^{-1}{\bf1}_{\{|v|<k\}}
       \begin{pmatrix}1&v\\2v&2v^2\end{pmatrix},
       \qquad v=(x-t)/s.
\]
Near \((0,S_0)\) these matrices are uniformly bounded.
Their change is uniformly small outside a band of arbitrarily
small width around \(\{-q,q\}\); the remaining empirical error
is bounded by a constant times that band's empirical mass.
The no-boundary-atom assumption and the strong law let this
bound tend to zero. Thus the empirical Jacobian averaged along
the segment from the true to the estimated pair tends almost
surely to its population value
\[
 D=-S_0^{-1}
       \begin{pmatrix}A&0\\0&2C\end{pmatrix},\quad
 A=F((-q,q))>0,\quad
 C=E_F[(X/S_0)^2{\bf1}_{\{|X|<q\}}]>0.
\]
The cross entries vanish by symmetry; positivity follows from
the normal component. The segment integral is valid for the
piecewise absolutely continuous score functions.
The exact estimating-equation expansion is
\[
 0=P_nm(\,\cdot\,;0,S_0)+D_n(T_n,S_n-S_0)^{\mathsf T},
 \qquad D_n\to D .
\]
Both centered empirical moments are \(O_P(n^{-1/2})\), by their
bounded variances and Chebyshev's inequality. Invertibility of
\(D\) makes the estimated vector \(O_P(n^{-1/2})\).
The first row of the expansion, with its limiting zero cross
entry, now gives
\[
 \sqrt n T_n=\frac{S_0}{A}\,
       n^{-1/2}\sum_i\psi(k,X_i/S_0)+o_P(1).
\]
The scalar i.i.d.\ central limit theorem yields variance
\(S_0^2\beta/A^2=E_F\psi(q,X)^2/A^2\).
No unbounded observation moment was used.

Uniform moment convergence and the same vanishing boundary-band
argument give convergence of the numerator average to \(\beta\)
and of the interior proportion to \(A\). Substitution proves
\(n\widehat v_n\to S_0^2\beta/A^2\).
When every residual is interior, the first equation gives the
arithmetic mean; cancelling \(S_n\) in the displayed variance
estimate gives the claimed classical expression.
\end{proof}
\end{annexitem}

\begin{annexitem}{annex:p15.huber.joint.variance.bound}{thm:p15-huber-joint-variance-bound}{The calculation behind historical Table IV}
\begin{proof}
The population second equation and \(q=kS_0\) imply
\[
 q^2r=(1-\varepsilon)B(q)+\varepsilon E_H\psi(q,X)^2
       \leq(1-\varepsilon)B(q)+\varepsilon q^2 .
\]
Thus \(B(q)/q^2\geq(r-\varepsilon)/(1-\varepsilon)\).
The function \(B(q)/q^2=E_\Phi\min(X^2/q^2,1)\) decreases
strictly and continuously from one to zero. For \(0\leq\varepsilon<r<1\)
the right side belongs to \((0,1)\), so it determines exactly one
positive \(q_{\max}\), and \(q\leq q_{\max}\).
Also \(A(q)/q=2q^{-1}\int_0^q\phi(x)\,dx\) decreases strictly,
because the normal density decreases on the positive half-line.
In the regular variance formula the numerator is \(q^2r\), while
the denominator's unsquared factor is
\[
 (1-\varepsilon)A(q)+\varepsilon H((-q,q))
       \geq(1-\varepsilon)A(q).
\]
After division by \(q\), its lower bound is at least
\((1-\varepsilon)A(q_{\max})/q_{\max}\).
This proves the variance inequality.
Symmetric contamination strictly outside the final cutoff makes
both inequalities equalities, has no atom at that cutoff, and
therefore lies in the regular domain attaining the bound.
The identity with \(W\) follows from the equation for \(q_{\max}\).
If \(\varepsilon\geq r\), that equation has a nonpositive right
side, whereas \(B(q)/q^2>0\) at every finite \(q\); a finite
solution cannot be asserted.
\end{proof}
\end{annexitem}

\begin{annexitem}{annex:p15.huber.mr.projection}{thm:p15-huber-mr-projection}{The needed projection and smooth-root theorem}
\begin{proof}
Conditional Jensen gives \(Ep(X)^2\leq Eh_0^2\) and \(Ep(X)=0\).
Let \(h^\circ(x_1,\ldots,x_r)=h_0(x_1,\ldots,x_r)-\sum_jp(x_j)\).
Its mean is zero, and its conditional mean given any one coordinate
is zero. In the average over subsets, each observation belongs to
\(\binom{n-1}{r-1}\) subsets. Consequently the sum of the projected
terms is exactly \(r n^{-1}\sum_i p(X_i)\).
This leaves \(R_n\), the subset average of \(h^\circ\).
Two such centered kernels on disjoint subsets are independent.
If they share exactly one index, conditional independence given
that observation and their zero conditional means make their
covariance zero. The number of ordered subset pairs sharing at
least two indices is \(O(n^{2r-2})\), for fixed \(r\);
each covariance is bounded in absolute value by \(E(h^\circ)^2\).
Division by \(\binom nr^2\) proves \(ER_n^2=O(n^{-2})\).
For \(r=1\) the remainder is identically zero.
The scalar i.i.d.\ central limit theorem and
\(\sqrt nR_n\to0\) in \(L^2\) prove the projection limit.

A subset average of any square-integrable kernel converges in
probability to its mean: only overlapping subsets can have
nonzero covariance, there are \(O(n^{2r-1})\) such pairs, and
the same variance bound is \(O(n^{-1})\).
Apply this fact to \(\partial_t h_0\).
The fundamental theorem of calculus and the uniform modulus
give
\[
 0=U_n(0)+T_nD_n,\qquad
 D_n=\binom nr^{-1}\sum_{|I|=r}
                 \int_0^1\partial_t h_{vT_n}(X_I)\,dv
       \longrightarrow -D
\]
in probability, since its difference from the derivative average
at zero is at most \(\omega(|T_n|)\) on the event where \(T_n\)
is sufficiently small. Dividing by the nonzero limiting denominator
proves the root limit and its variance.

For the final bridge, the chain rule gives
\(-\partial_t h_t|_0=\sum_{j=1}^r\partial_j h_0\).
Symmetry and absolute integrability make its expectation
\(rE\partial_1h_0\). The specified compact-interval dominating
bound permits differentiating the conditional integral for \(p(x)\)
by dominated convergence, giving
\(p'(x)=E[\partial_1h_0(x,X_2,\ldots,X_r)]\).
Fubini now gives \(Ep'(X)=E\partial_1h_0\), proving the factor \(r\).
This states the actual derivative and interchange conditions behind
the abbreviated source variance.
\end{proof}
\end{annexitem}

\begin{annexitem}{annex:p15.huber.pair.median}{thm:p15-huber-pair-median}{The distinct pair-average example and its comparisons}
\begin{proof}
For each distinct pair,
\(|(X_i-t)+(X_j-t)|=2|Z_{ij}-t|\).
The ordinary finite-sample median minimization already proved
therefore identifies the minimizers exactly.
For the asymptotics define
\[
 H(t)=P((X_1+X_2)/2\leq t),\qquad
 U_n(t)=\binom n2^{-1}\sum_{i<j}{\bf1}_{\{Z_{ij}\leq t\}}.
\]
The density of this pair average is
\(2\int f(x)f(2t-x)\,dx\). It is continuous at zero and equals
\(2\int f^2\) there. Here is the needed continuity justification:
translation is continuous in \(L^2(dx)\), first for continuous
compactly supported functions by uniform continuity and bounded
support, then for every \(L^2\) function by approximation with
such functions and the translation invariance of its norm.
The approximation follows by truncating tails, approximating
simple functions on a finite interval, and continuous ramps.
Cauchy--Schwarz now makes the displayed convolution continuous.
Symmetry gives \(H(0)=1/2\), and hence
\[
 H(t)=1/2+2t\int f^2+o(|t|).
\]

For the bounded indicator kernel, its first centered projection is
\(p_t(x)=F(2t-x)-H(t)\).
At zero, symmetry and continuity of \(F\) give
\(p_0(X)=1/2-F(X)\).
The random variable \(F(X)\) is uniform on \((0,1)\):
continuity of the CDF supplies a point at each intermediate
probability value, and any plateau has zero probability.
Thus \(Ep_0^2=1/12\).
For \(t_n=u/\sqrt n\), bounded convergence gives
\(p_{t_n}\to p_0\) in \(L^2(F)\).
The centered projection sum has a bounded triangular-array
Lindeberg limit with variance \(4(1/12)=1/3\).
The remainder argument of the preceding proof is uniform in
\(t_n\), because all indicator kernels are bounded by one;
its variance is \(O(n^{-2})\). Therefore
\[
 \sqrt n\{U_n(t_n)-H(t_n)\}\Longrightarrow N(0,1/3).
\]
At a fixed small positive or negative \(t\), the same variance
bound gives convergence of \(U_n(t)\) to \(H(t)\), proving
consistency of its median.
For the limit, the two central order-statistic thresholds of
the \(\binom n2\) pair values differ from \(1/2\) by
\(O(\binom n2^{-1})=o(n^{-1/2})\).
The midpoint lies between those two order statistics.
Combining their indicator-count bounds with the last normal limit
and the expansion of \(H(t_n)\) gives variance
\[
 \frac{1/3}{(2\int f^2)^2}
             =\frac1{12(\int f^2)^2}.
\]
This argument never treats the overlapping pairs as independent.

Write \(a=1-\varepsilon\). For density contamination,
\[
 \int f^2=a^2\int\phi^2+2a\varepsilon\int\phi h
                      +\varepsilon^2\int h^2
       \geq\frac{a^2}{2\sqrt\pi}.
\]
For \(0<\varepsilon<1\), the cross integral is strictly positive
because \(\phi\) is positive and \(h\) is a probability density.
This gives the strict individual bound and the equality at zero
contamination. Taking symmetric uniform contaminants on
\([-M,M]\) makes \(\int h^2=1/(2M)\) and
\(\int\phi h\leq1/(2M)\), so the bound is approached.
This proves its sharp supremum.

For two centered normal densities with standard deviations
\(s,t>0\), multiplying their exponentials and integrating gives
\[
 \int\phi_s(x)\phi_t(x)\,dx
            =\{2\pi(s^2+t^2)\}^{-1/2}.
\]
Substituting \(s,t\in\{1,3\}\) and expanding the mixture square
proves its stated exact integral.
For the Proposal 2 comparison,
\(E_{N(0,9)}\psi(q,X)^2=9B(q/3)\) by \(X=3Z\), and
\(P_{N(0,9)}(|X|<q)=A(q/3)\).
Substitution into its population scale equation and the proved
regular variance formula gives both expressions in the body.
The source's rounded comparison rows are evaluations of these
formulas.
\end{proof}
\end{annexitem}
